\chapter{Derivation of the Damped Random Walk Transition Probability}
\label{app:transition}

The transition probability of Equation~\ref{eq:drw_transition} follows from solving the Ornstein-Uhlenbeck stochastic differential equation. Here $X_t$ denotes the process $f(t)$ of Equation~\ref{eq:drw_sde}, rewritten using the standard It\^o notation $\mathrm{d}X_t$ and the damping rate $\theta = 1/\tau$:
\begin{equation}
    \label{eq:app_ou}
    \mathrm{d}X_t = \theta\,(\mu - X_t)\,\mathrm{d}t + \sigma\sqrt{2\theta}\,\mathrm{d}W_t,
\end{equation}
where $\mu$ is the long-term mean, $\theta$ the mean-reversion rate, $\sigma$ the stationary standard deviation, and $W_t$ a Wiener process. The equation is solved by variation of parameters, using the integrating factor $e^{\theta t}$. Defining $Y_t = X_t\,e^{\theta t}$ and applying It\^o's lemma,
\begin{equation}
    \label{eq:app_ito}
    \mathrm{d}Y_t = \theta X_t\, e^{\theta t}\,\mathrm{d}t + e^{\theta t}\,\mathrm{d}X_t
    = e^{\theta t}\,\theta\mu\,\mathrm{d}t + \sigma\sqrt{2\theta}\, e^{\theta t}\,\mathrm{d}W_t,
\end{equation}
where the drift terms in $X_t$ cancel on substituting Equation~\ref{eq:app_ou}. Integrating from $0$ to $t$,
\begin{equation}
    \label{eq:app_integrated}
    X_t\, e^{\theta t} - X_0
    = \theta\mu \int_0^t e^{\theta s}\,\mathrm{d}s
    + \sigma\sqrt{2\theta} \int_0^t e^{\theta s}\,\mathrm{d}W_s
    = \mu\left(e^{\theta t} - 1\right) + \sigma\sqrt{2\theta} \int_0^t e^{\theta s}\,\mathrm{d}W_s,
\end{equation}
and multiplying through by $e^{-\theta t}$ gives the solution
\begin{equation}
    \label{eq:app_solution}
    X_t = X_0\, e^{-\theta t} + \mu\left(1 - e^{-\theta t}\right)
    + \sigma\sqrt{2\theta} \int_0^t e^{-\theta(t-s)}\,\mathrm{d}W_s.
\end{equation}

The first two terms are deterministic, since $X_0$, $\mu$, and $\theta$ are fixed. The final term is an It\^o integral of a deterministic function against the Wiener process, and is therefore Gaussian: it can be approximated as a weighted sum of independent Gaussian increments $\mathrm{d}W_s \sim \sqrt{\mathrm{d}s}\,\varepsilon_s$ with $\varepsilon_s \sim \mathcal{N}(0,1)$, and any linear combination of Gaussians is itself Gaussian. The state $X_t$ conditioned on $X_0$ is thus normally distributed, and it remains to compute its mean and variance.

The mean is obtained by taking the expectation of Equation~\ref{eq:app_solution}. The It\^o integral has zero mean, so
\begin{equation}
    \label{eq:app_mean}
    \mathbb{E}\!\left[X_t \mid X_0\right]
    = X_0\, e^{-\theta t} + \mu\left(1 - e^{-\theta t}\right).
\end{equation}
The variance comes from the It\^o isometry, which turns the variance of the stochastic integral into an ordinary integral of the squared integrand,
\begin{equation}
    \label{eq:app_var}
    \mathrm{Var}\!\left[X_t \mid X_0\right]
    = 2\theta\sigma^2 \int_0^t e^{-2\theta(t-s)}\,\mathrm{d}s
    = 2\theta\sigma^2 \cdot \frac{1 - e^{-2\theta t}}{2\theta}
    = \sigma^2\left(1 - e^{-2\theta t}\right).
\end{equation}

Collecting the mean and variance, the transition probability is
\begin{equation}
    \label{eq:app_transition}
    X_t \mid X_0 \;\sim\; \mathcal{N}\!\left(\mu + \left(X_0 - \mu\right)e^{-\theta t},\;\; \sigma^2\left(1 - e^{-2\theta t}\right)\right).
\end{equation}
Writing the elapsed time as $\Delta t$ and the damping timescale as $\tau = 1/\theta$, this is Equation~\ref{eq:drw_transition}. The factor $\sqrt{2\theta}$ in the driving term of Equation~\ref{eq:app_ou} is what makes $\sigma$ the stationary standard deviation directly: as $t \to \infty$ the variance tends to $\sigma^2$, independently of $\tau$.


\chapter{Neural Stochastic Flow in Depth}
\label{app:loss}

This appendix supplements the description of the reconstruction method in Chapter \ref{chap:nsf}, following \cite{Kiyohara2025SNF}. It first gives the architecture of the bridge model, the counterpart of the flow whose construction is set out in the main text, and then makes explicit the triplet-consistency loss through which the two models are trained jointly.

\section{The Bridge Architecture}
\label{app:bridge_arch}

The flow $F_\theta$ of Chapter \ref{chap:nsf} models the transition between two states separated by an elapsed time, and its construction, a state-dependent Gaussian base distribution followed by conditioned affine-coupling layers, is given in Equations \ref{eq:base} and \ref{eq:coupling}. The bridge $q_\xi$ of Equation \ref{eq:bridge} is built in the same way, but it models a different conditional where the intermediate state $x_j$ given \textit{both} bracketing states $x_i$ and $x_k$ and the three times, rather than a future state given a single past one. It is this conditioning on both endpoints that makes the bridge suitable for gap filling, where the observations on either side of a gap are known.

The bridge conditioner is
\begin{equation}
    \label{eq:app_bridge_cond}
    c_b := \big(x_i,\, x_k,\, \Delta t_{ij},\, \Delta t_{jk}\big),
\end{equation}
carrying both surrounding states and the two elapsed times to the intermediate point, in place of the single pair $(x_s, \Delta t)$ used by the flow. As with the flow, a sample of the intermediate state is drawn by first initializing a Gaussian base distribution and then transforming it through a sequence of $L$ conditioned affine-coupling layers \parencite{Dihn2015Coupling},
\begin{equation}
    \label{eq:app_bridge_sample}
    x_j = G_\xi(z_b, c_b) = g_L(\,\cdot\,; c_b, \xi_L) \circ \cdots \circ g_1(z_b; c_b, \xi_1),
    \qquad z_b \sim \mathcal{N}\big(\mu_b(c_b),\, \sigma_b^2(c_b)\big),
\end{equation}
where the mean and variance of the base distribution and the scale and shift of each coupling layer are produced by multilayer perceptrons taking the conditioner $c_b$ as input, in the same manner as Equations \ref{eq:base} and \ref{eq:coupling} for the flow. The bridge and the flow therefore share the same affine-coupling structure and differ only in what they are conditioned on, so that both admit a tractable likelihood through the change-of-variables formula of Equation \ref{eq:changeofvars} and can be sampled in a single pass.

\section{The Triplet-Consistency Loss}
\label{app:triplet_loss}

The training objective of Equation \ref{eq:loss} combines a data term and a flow-consistency term, both evaluated on the triplets $(t_i, t_j, t_k)$ sampled from each light curve. This section writes out both terms explicitly for a single triplet, the data term as a sum of transition log-densities and the flow-consistency term as a difference of log-densities computed on Monte Carlo samples drawn from the models.

\subsection{Data Term}

The data term is the negative log-likelihood of the transitions of the triplet under the flow. All three ordered pairs are included, the two one-step transitions and the direct endpoint transition, so that
\begin{equation}
    \label{eq:app_ldata}
    \ell_\mathrm{data} = -\frac{1}{3}\Big[
        \log p_\theta(x_j \mid x_i, \Delta t_{ij})
        + \log p_\theta(x_k \mid x_j, \Delta t_{jk})
        + \log p_\theta(x_k \mid x_i, \Delta t_{ik})
    \Big],
\end{equation}
each log-density being evaluated through the change-of-variables formula of Equation \ref{eq:changeofvars}. Including the endpoint transition $x_i \rightarrow x_k$ alongside the two one-step transitions is what trains the flow on the full span of the triplet directly, rather than leaving the long transition to be reproduced by composing shorter ones. The data term of the objective is the expectation of $\ell_\mathrm{data}$ over the triplets.

\subsection{Flow-Consistency Term}

The flow-consistency term of Chapter \ref{chap:nsf} is written in Equations \ref{eq:flow12} and \ref{eq:flow21} as two variational bounds on the Kullback-Leibler divergence between the one-step and two-step transitions. What follows makes explicit how those bounds are evaluated in practice, as a difference of log-densities computed on Monte Carlo samples, for a single triplet.

Consider one light curve and three times $t_i < t_j < t_k$, with corresponding states $x_i$, $x_j$, and $x_k$. The flow defines the one-step transition densities
\begin{equation}
    \label{eq:app_onestep}
    p_\theta(x_k \mid x_i, \Delta t_{ik}), \qquad
    p_\theta(x_j \mid x_i, \Delta t_{ij}), \qquad
    p_\theta(x_k \mid x_j, \Delta t_{jk}),
\end{equation}
and the bridge defines the intermediate conditional $q_\xi(x_j \mid x_i, x_k)$. The joint distribution over the intermediate and final states $(x_j, x_k)$ can be written in two ways, and the loss enforces that they assign consistent log-densities to samples generated from the models.

\subsubsection{Forward Direction}

In the forward direction the final state is drawn from the one-step flow and the intermediate state from the bridge,
\begin{equation}
    \label{eq:app_fwd_sample}
    x_k \sim p_\theta(\,\cdot \mid x_i, \Delta t_{ik}), \qquad
    x_j \sim q_\xi(\,\cdot \mid x_i, x_k),
\end{equation}
which corresponds to the factorization
\begin{equation}
    \label{eq:app_fwd_fact}
    p_\theta(x_k \mid x_i, \Delta t_{ik})\; q_\xi(x_j \mid x_i, x_k).
\end{equation}
The same joint distribution can be written as a two-step flow through the intermediate state,
\begin{equation}
    \label{eq:app_fwd_twostep}
    p_\theta(x_j \mid x_i, \Delta t_{ij})\; p_\theta(x_k \mid x_j, \Delta t_{jk}).
\end{equation}
Given the sample $(x_j, x_k)$, the contribution to the loss is the difference of the log-densities of the two factorizations,
\begin{equation}
    \label{eq:app_l12}
    \ell_{1\text{-}\mathrm{to}\text{-}2}
    = \log p_\theta(x_k \mid x_i, \Delta t_{ik})
    + \log q_\xi(x_j \mid x_i, x_k)
    - \log p_\theta(x_j \mid x_i, \Delta t_{ij})
    - \log p_\theta(x_k \mid x_j, \Delta t_{jk}).
\end{equation}

\subsubsection{Reverse Direction}

In the reverse direction both states are generated by two consecutive flow transitions,
\begin{equation}
    \label{eq:app_rev_sample}
    x_j \sim p_\theta(\,\cdot \mid x_i, \Delta t_{ij}), \qquad
    x_k \sim p_\theta(\,\cdot \mid x_j, \Delta t_{jk}),
\end{equation}
and this is compared with the factorization that uses the bridge and the direct one-step flow, $q_\xi(x_j \mid x_i, x_k)\,p_\theta(x_k \mid x_i, \Delta t_{ik})$. The corresponding log-density difference is
\begin{equation}
    \label{eq:app_l21}
    \ell_{2\text{-}\mathrm{to}\text{-}1}
    = \log p_\theta(x_j \mid x_i, \Delta t_{ij})
    + \log p_\theta(x_k \mid x_j, \Delta t_{jk})
    - \log q_\xi(x_j \mid x_i, x_k)
    - \log p_\theta(x_k \mid x_i, \Delta t_{ik}).
\end{equation}

\subsubsection{Combined Loss}

The total consistency loss for the triplet is the sum of the two directions,
\begin{equation}
    \label{eq:app_ltriplet}
    \ell_\mathrm{flow} = \ell_{1\text{-}\mathrm{to}\text{-}2} + \ell_{2\text{-}\mathrm{to}\text{-}1},
\end{equation}
and the flow-consistency term of the objective is its expectation over triplets and over the Monte Carlo samples drawn from the models. Combined with the data term of Equation \ref{eq:app_ldata}, the full per-triplet loss is
\begin{equation}
    \label{eq:app_full}
    \ell = \ell_\mathrm{data} + \lambda\,\ell_\mathrm{flow},
\end{equation}
matching the objective of Equation \ref{eq:loss}. When the flow and bridge are mutually consistent, the two factorizations of the joint distribution over $(x_j, x_k)$ coincide, and the expected log-density differences of Equations \ref{eq:app_l12} and \ref{eq:app_l21} vanish.